\chapter{The Machine's Inputs: How the Eyes, Ears, and Other Sensors Connect to It}
\chaptermark{The Machine's Inputs}
\label{sec:sensors}

\section{A sensor as a transducer}

In Fig.~\ref{fig:machine} data entered the machine from the left, from the ``sensors'' block. Now we open that block. To an engineer every sense organ is a \emph{transducer} with an analog front end and an analog-to-digital converter at the end, except that the digital output is not numbers but spikes.

All sense organs share the same scheme (Fig.~\ref{fig:transducer}). A physical quantity (light, pressure, vibration, a molecule) acts on a receptor protein in the membrane of a sensory cell. This protein works as a gate: by itself or through a short chain of reactions, it opens an ion channel. A \emph{receptor current} arises, and with it an analog voltage across the membrane. The voltage then passes through gain control and enters a spike encoder, the familiar integrating neuron~\eqref{eq:glutneuron}, which turns a level into a rate and spike times. The output is almost always the same: the sensory neurons of the eye, the ear, smell, and skin release glutamate, so the inputs connect directly to the \nt{GLUT} bus.

\begin{figure}[t]
\centering
\begin{tikzpicture}[>=Stealth,
  blk/.style={draw,very thick,rounded corners=2pt,align=center,font=\footnotesize,minimum height=1.0cm,text width=1.9cm,fill=blue!5}]
\node[align=center,font=\small] (P) at (0.1,0) {light, sound,\\pressure,\\molecule};
\node[blk] (R) at (3.0,0) {receptor\\gate};
\node[blk] (G) at (6.5,0) {gain\\control};
\node[blk] (E) at (10.0,0) {spike\\encoder};
\node[align=center,font=\small] (B) at (13.4,0) {\nt{GLUT}\\bus};
\draw[->,very thick] (P) -- (R);
\foreach \ic/\xx in {sun/-1.1,speaker/-0.3,press/0.5,molecule/1.3}{\pic[scale=0.85] at (\xx,1.05) {\ic};}
\draw[->,very thick] (R) -- node[above,font=\scriptsize] {current} (G);
\draw[->,very thick] (G) -- node[above,font=\scriptsize] {level} (E);
\draw[->,very thick,red!70!black] (E) -- node[above,font=\scriptsize,black] {spikes} (B);
\draw[decorate,decoration={brace,amplitude=5pt,mirror},gray!60!black] (1.8,-0.8) -- (7.7,-0.8) node[midway,below=5pt,font=\scriptsize,black] {analog part};
\draw[decorate,decoration={brace,amplitude=5pt,mirror},gray!60!black] (8.8,-0.8) -- (11.2,-0.8) node[midway,below=5pt,font=\scriptsize,black] {spikes};
\end{tikzpicture}
\caption{Any sense organ as a chain of conversions. The receptor protein opens a channel and produces a current; gain control adapts it to the conditions; the encoder~\eqref{eq:glutneuron} turns the level into spikes, which go out onto the \nt{GLUT} bus. In the eye and the ear the first stages produce no spikes at all: the signal stays analog until it has to be sent far.}
\label{fig:transducer}
\end{figure}

One detail is very familiar to an electronics engineer. In the eye and the ear the very first cells produce no spikes at all: they pass a smooth voltage to their neighbors, like an analog stage on a circuit board. Spikes appear only where the signal must travel far. Over millimeters an analog signal decays and picks up noise, while a spike, as we saw in Chapter~\ref{sec:neurontypes}, reaches the end of the wire at full height. Digitization happens where the long line begins.

\section{At the limit of physics}

The brain's sensors work at the physical limits of sensitivity. A light sensor in the dark responds to a \emph{single photon}: absorbing one particle of light gives a current of about a picoampere, visible above the sensor's own noise~\cite{Baylor1979}. A sound sensor detects a displacement of its sensitive bundle of less than a nanometer~\cite{Hudspeth1989}.

When light is scarce, precision is limited not by the sensor but by the light itself: photons arrive at random, as a Poisson stream. If on average $N=\Phi T$ photons reach the sensor during the integration time $T$, their count fluctuates by $\sqrt N$. For an increment $\Delta N$ to be noticeable, it must be several times larger than this fluctuation, and the detectable fractional change in brightness is
\begin{empheq}[box=\keybox]{equation}
\frac{\Delta I}{I} \;\approx\; \frac{k}{\sqrt{\Phi\,T}} .
\label{eq:photonnoise}
\end{empheq}
This is the same formula as for the spike count~\eqref{eq:rateerror}: the shot noise of photons and the shot noise of spikes obey one law. In the dark the eye has only one way to improve precision, which is to collect photons longer, and the integration time does indeed grow to tenths of a second. That is why we see more slowly at dusk.

\section{Dynamic range: gain control}

The hardest engineering problem for sensors is range. Illuminance varies by about ten orders of magnitude from a starlit night to sunlit snow, and sound intensity by twelve from the threshold of hearing to the threshold of pain. Yet the spike rate varies from zero to a few hundred per second, less than three orders of magnitude. Such an input cannot be mapped linearly onto such an output.

The solution is the same as in any receiver: \emph{automatic gain control}. The responses of the sensory cells of the eye are well described by the formula
\begin{empheq}[box=\keybox]{equation}
r \;=\; r_{\max}\,\frac{I}{I+\sigma} ,
\label{eq:nakarushton}
\end{empheq}
where $\sigma$ is the brightness at which the response reaches half its maximum~\cite{NakaRushton1966}. By itself,~\eqref{eq:nakarushton} covers only two or three orders of magnitude. But $\sigma$ is not constant: during prolonged work against a bright background it grows along with the mean brightness. The response curve shifts along the logarithmic brightness axis and each time places its steep part where the input currently is (Fig.~\ref{fig:adaptation}). This is the same division by total activity as the shunting inhibition of Chapter~\ref{sec:gaba}~\cite{Carandini2012}, except that the denominator here is the mean brightness over the last few seconds.

\begin{figure}[t]
\centering
\begin{tikzpicture}[>=Stealth,x=1.25cm,y=2.8cm]
\draw[->,gray!70!black] (-2,0) -- (6.4,0) node[below left,font=\small,black] {$\log_{10} I$};
\draw[->,gray!70!black] (-2,0) -- (-2,1.15) node[left,font=\small,black] {$r/r_{\max}$};
\foreach \u in {-2,0,2,4,6}{\draw[gray!60!black] (\u,-0.02) -- (\u,0.02); \node[below,font=\scriptsize] at (\u,0) {$\u$};}
\draw[densely dashed,gray!60!black] (-2,0.5) -- (6.2,0.5);
\foreach \s/\c/\lab in {0/blue!60!black/{dark},2/green!45!black/{medium},4/red!70!black/{bright}}{
  \pgfmathsetmacro{\lo}{max(-2,\s-3)}
  \draw[very thick,\c] (-2,0) -- (\lo,0) plot[domain=\lo:6.2,samples=80] (\x,{1/(1+10^(\s-\x))});
  \node[font=\scriptsize,\c,anchor=west] at (\s+0.15,0.42) {\lab};}
\foreach \s/\ic in {0/moon,2/cloud,4/sun}{\pic at (\s+0.6,0.64) {\ic};}
\end{tikzpicture}
\caption{Gain control of a light sensor by~\eqref{eq:nakarushton}. Each curve covers two or three orders of magnitude of brightness; on moving to a brighter background $\sigma$ grows, and the curve shifts along the logarithmic axis. Together the shifting curves cover the whole range.}
\label{fig:adaptation}
\end{figure}

Gain control has a price, familiar from Chapter~\ref{sec:code}: the sensor reports not brightness but brightness \emph{relative to the background}. A constant input gradually stops evoking a response, while every change evokes one. From the very start, the machine's inputs speak of changes, not of levels, and this is the same economy of spikes as in Proposition~\ref{prop:economy}~\cite{Barlow1961}. Hence also the on and off sensors of Chapter~\ref{sec:glutcomputer}~\cite{Schiller1992}: some report that it got brighter, others that it got darker.

Engineers have copied this trick in \emph{event cameras}. Such a camera does not take frames on a timer: each of its pixels, on its own and with no shared clock, emits a ``brighter'' or ``darker'' event when the logarithm of brightness has changed by a set fraction~\cite{Lichtsteiner2008}. This is exactly the two-wire on/off code of Chapter~\ref{sec:glutcomputer}, implemented in silicon, and it gives the camera a range and speed out of reach of an ordinary sensor array~\cite{Gallego2022}.

\section{The eye: compression down to a cable}

The eye has about $10^8$ light sensors~\cite{Curcio1990}, while the output neurons that carry the image to the brain number about a million. Before the signal leaves the eye, it is compressed about a hundredfold. The main compression tricks are familiar. Each output neuron responds to the difference between the brightness in a small spot and the brightness around it; this is subtraction of neighbors through inhibition, as in Chapter~\ref{sec:gaba}. Uniform regions of the picture cost almost nothing, while edges and changes are transmitted. Different output neurons are specialized: some report brightening, others darkening, still others motion in a particular direction (Fig.~\ref{fig:direction}); in the mouse more than thirty such output channels have been counted~\cite{Baden2016}.

What is the capacity of the resulting cable? In the guinea pig, recordings from output neurons gave about $875$ thousand bits per second for $\sim10^5$ such neurons; scaling to the human million output neurons gives on the order of $10^7$ bits per second~\cite{Koch2006}, as much as an old ten-megabit network cable. At a mean rate of a few spikes per second per neuron this is a few bits per spike, as we found in Chapter~\ref{sec:code}.

\section{The ear: a filter bank}

The ear solves a different problem: sound must be decomposed by frequency. An engineer would install a bank of bandpass filters, and the ear does exactly that, mechanically. The sound vibration travels along an elastic partition whose properties change smoothly along its length, and each section resonates at its own frequency. The sensors sitting on a section respond only to their own band (Fig.~\ref{fig:filterbank}). Frequencies are laid out by place roughly logarithmically: each octave gets a similar share of the sensors. The index of the sensor that fired is the frequency, the place code of Chapter~\ref{sec:code}.

\begin{figure}[t]
\centering
\begin{tikzpicture}[>=Stealth,x=1.35cm,y=2.2cm]
\draw[->,gray!70!black] (-0.3,0) -- (7.6,0) node[above left,font=\small,black] {frequency, Hz};
\draw[->,gray!70!black] (-0.3,0) -- (-0.3,1.15) node[left,font=\small,black] {response};
\foreach \k/\f in {0/125,1/250,2/500,3/1000,4/2000,5/4000,6/8000}{
  \draw[gray!60!black] (\k,-0.02) -- (\k,0.02); \node[below,font=\scriptsize] at (\k,0) {\f};
  \draw[thick,blue!60!black] plot[domain={max(\k-1.2,-0.3)}:{\k+0.7},samples=60] (\x,{1/(1+((\x-\k)/(\x<\k?0.45:0.2))^4)});}
% a piano keyboard: seven white keys per octave, like the octaves on the axis
\foreach \i in {0,...,48}{\draw[gray!50!black,fill=white,line width=0.3pt] ({\i/7},-0.44) rectangle ({(\i+1)/7},-0.26);}
\foreach \o in {0,...,6}{\foreach \j in {0,1,3,4,5}{\fill[black] ({\o+(\j+1)/7-0.035},-0.35) rectangle ({\o+(\j+1)/7+0.035},-0.26);}}
\end{tikzpicture}
\caption{The ear as a bank of bandpass filters, schematic. Each curve is the response of the sensors of one section to sounds of different frequency; the center frequencies are an octave apart, like keys on a logarithmic scale. Real filters have a steep high-frequency slope and a gentle low-frequency one. Below is a keyboard: on it, as in the ear, each octave gets the same amount of space.}
\label{fig:filterbank}
\end{figure}

The resonators of the ear are not passive. Some of the sensors themselves drive the partition in time with the sound and amplify weak vibrations thousands of times in amplitude (by $66$--$76$\,dB), and the gain falls as loudness grows~\cite{Ruggero1997,Robles2001}. To an engineer this is an amplifier with positive feedback, set right at the edge of self-oscillation: the same life on the edge as in the network of Chapter~\ref{sec:gaba}, where near the boundary the system is especially sensitive to small signals. Loudness compression is done by the same amplifier: quiet sounds are amplified strongly, loud ones weakly.

The ear also has a second code. At low frequencies neurons fire in phase with the sound: a neuron fires at a particular part of each wave, though not on every wave. In the guinea pig, phase locking begins to weaken at about $600$\,Hz and is still noticeable up to about $3.5$\,kHz~\cite{PalmerRussell1986}. In this way spike times preserve the precise timing of the sound, and from it the difference in arrival time at the two ears, from which the network of Chapter~\ref{sec:glutcomputer} finds the direction~\cite{Grothe2010}. At high frequencies the spikes cannot keep up with the wave, and only the place code remains. A single ear uses both codes of Chapter~\ref{sec:code}: timing where the neurons can keep up, and place where they cannot.

\section{The other sensors}

\begin{table}[t]
\centering
\footnotesize
\setlength{\tabcolsep}{3pt}
\renewcommand{\arraystretch}{1.15}
\begin{tabular}{>{\raggedright}p{1.9cm}>{\raggedright}p{2.6cm}>{\raggedright}p{4.0cm}>{\raggedright\arraybackslash}p{4.6cm}}
\toprule
Sense & What it measures & How the input is built & Trick from the book \\
\midrule
vision & light & $\sim10^8$ sensors, compression by $\sim100$, $\sim10^7$ bits/s & gain control, subtraction of neighbors, two wires, direction of motion \\
hearing & air pressure & mechanical filter bank, active amplifier & place code and timing code, delays \\
touch & pressure, vibration & fast- and slow-adapting sensors & a ``high-pass filter --- low-pass filter'' pair \\
temperature & heat & separate sensors for warm and cold & two-wire code \\
smell & molecules & $\sim400$ receptor types, one type per sensor & combinatorial code: an odor is the set of types that fired \\
taste & substances in food & five qualities: sweet, bitter, salty, sour, umami & labeled lines; some cells release ATP or serotonin rather than glutamate \\
body position & muscle stretch & sensors of stretch length and velocity & feedback for the movement controller \\
\bottomrule
\end{tabular}
\caption{How the sensors are connected. All the devices in the third column have been measured; the right column links them to tricks from earlier chapters.}
\label{tab:sensors}
\end{table}

The other senses are collected in Table~\ref{tab:sensors}, and almost every row shows a trick from earlier chapters. Touch uses a pair of filters: some pressure sensors respond to contact and fall silent at once, while others keep responding as long as the pressure lasts. Temperature is carried by two wires, separate for warm and cold. The most unusual input is smell. Humans have about four hundred types of olfactory receptor, and each olfactory neuron carries only one type~\cite{Mombaerts2004}. An odor activates not one type but a set, and the message is \emph{which} types fired. To an engineer this is a binary vector of length about four hundred, with an astronomical number of possible values.

\begin{keyblock}
\begin{proposition}[The machine's inputs]
\label{prop:sensors}
Each sense is a transducer that works at the physical limit of sensitivity, compresses a huge dynamic range by automatic gain control, and sends mainly \emph{changes} over the \nt{GLUT} bus. To do this, the sensors use the same tricks as the machine itself: the two-wire code, the place code, the timing code, and division by the background. The structure of the sensors has been measured; that their codes are tuned for the greatest information per spike is a well-founded hypothesis.
\end{proposition}
\end{keyblock}

The inputs, like everything else, work asynchronously. Each sensor emits a spike when it has something to report, and the different senses arrive with different delays: sound after milliseconds, light after tens of milliseconds. How the machine combines them into one picture of the world with no shared clock is one of the problems of coordination in time from Chapter~\ref{sec:synthesis}.

\section{Exercises}

\begin{exercise}[\lvlA]
\label{ex:11:1}
\emph{How long to collect photons.} At dusk $100$ photons per second reach a sensor; an increment is noticeable if it is $3$ times the fluctuation~\eqref{eq:photonnoise}. (a)~How long does one sensor need to notice a $10\,\%$ change in brightness? (b)~How many sensors must be summed to manage it within $0.2\,\text{s}$, and by what factor does the spatial resolution along one axis drop?
\end{exercise}

\begin{exercise}[\lvlA]
\label{ex:11:2}
\emph{How many bits per spike.} (a)~The eye transmits $\sim10^7$ bits/s over $10^6$ output neurons firing $3$--$5$ spikes per second. What fraction of the capacity~\eqref{eq:spikeentropy} at $\Delta t=1\ms$ does it use? (b)~An odor activates $20$ of $\approx400$ receptor types. How many bits does such a set carry, and how many types do two random odors share on average?
\end{exercise}

\begin{exercise}[\lvlB]
\label{ex:11:3}
\emph{What one curve~\eqref{eq:nakarushton} can do.} (a)~Over what range of brightness does the response rise from $10\,\%$ to $90\,\%$ of $r_{\max}$? How many orders of magnitude is that? (b)~How many positions of the shifting curve are needed to cover ten orders of magnitude of brightness? Twelve orders of magnitude of loudness?
\end{exercise}

\begin{exercise}[\lvlB]
\label{ex:11:4}
\emph{The limit of the timing code in the ear.} The difference in arrival time at the two ears is up to $0.22\,\text{m}/343\,\text{m/s}$. (a)~Spikes are phase-locked to a tone: up to what frequency do they determine the direction unambiguously? (b)~A spike jitters by $0.5\ms$, and $20$~\textmu s must be resolved. How many neurons at $100$ spikes per second must be averaged over $50\ms$?
\end{exercise}

\begin{exercise}[\lvlB]
\label{ex:11:5}
\emph{An event camera on the eye's cable.} A camera of $10^6$ pixels with an output of $10^7$ bits/s sends an event (address and sign) when $\ln I$ in a pixel has changed by $\ln1.2$. What is the highest speed at which an edge $500$ pixels long with a brightness step of $4$ can move? How many frames per second would an ordinary camera at $8$ bits per pixel get through the same output?
\end{exercise}

\begin{exercise}[\lvlB]
\label{ex:11:6}
\emph{Simulation: gain control.} The sensor~\eqref{eq:nakarushton} adjusts $\sigma$ with $\tau=1\,\text{s}$ either logarithmically, $\tau\,\dd\ln\sigma/\dd t=\ln I-\ln\sigma$, or linearly, $\tau\,\dd\sigma/\dd t=I-\sigma$; the background jumps $1\to100\to1$. For each law, how many seconds after the upward and the downward jump does the response to a $+10\,\%$ probe recover to half its adapted value?
\end{exercise}

\begin{exercise}[\lvlC]
\label{ex:11:7}
\emph{What world the curve is tuned for.} With small constant output noise, the information about brightness is greatest when $r(I)=r_{\max}\Phi_I(I)$, where $\Phi_I$ is the distribution function of brightness. (a)~For which distribution is the curve~\eqref{eq:nakarushton} optimal, and what is its spread in $\log_{10}I$? (b)~Which curve is optimal with Poisson output noise?
\end{exercise}
